% =============================================================================
\section{Transformada $Z$: ROC, inversa e estabilidade}\label{sec:z}
% =============================================================================

\subsection*{Teoria mínima}

\[
  X(z)=\sum_{n=-\infty}^{\infty}x[n]z^{-n},\qquad z\in\mathbb{C},\qquad
  x[n]=\frac{1}{2\pi j}\oint X(z)z^{n-1}dz .
\]
A transformada só está definida junto com sua \textbf{região de convergência
(ROC)}: anéis do plano $z$ delimitados pelos polos. Pares essenciais:
\[
  \delta[n]\leftrightarrow1\ (\text{todo }z),\qquad
  u[n]\leftrightarrow\frac{1}{1-z^{-1}},\ |z|>1,
\]
\[
  \alpha^nu[n]\leftrightarrow\frac{1}{1-\alpha z^{-1}},\ |z|>|\alpha|,
  \qquad
  -\alpha^nu[-n-1]\leftrightarrow\frac{1}{1-\alpha z^{-1}},\ |z|<|\alpha| .
\]
Note: \emph{a mesma expressão algébrica} corresponde a duas sequências
diferentes, dependendo da ROC. Propriedades úteis:
\[
  x[n-n_0]\leftrightarrow z^{-n_0}X(z),\qquad
  x_1[n]*x_2[n]\leftrightarrow X_1(z)X_2(z).
\]

\begin{teo}[causalidade e estabilidade no plano $z$]
Para um sistema causal, a ROC é $|z|>|p|_{\max}$. O sistema é
\textbf{BIBO estável} $\iff$ a ROC contém o \textbf{círculo unitário}
$|z|=1$. Logo, um sistema causal é estável $\iff$ \textbf{todos os polos têm
$|p_i|<1$}.
\end{teo}

\paragraph{Transformada inversa por frações parciais.} Escreva
$X(z)=\sum_i\frac{A_i}{1-p_iz^{-1}}$ (polos simples). Então
\[
  x[n]=\sum_i A_i\,p_i^{\,n}u[n]\quad\text{(se ROC ``para fora'')}.
\]
Atalho prático para $X(z)$ de ordem igual no numerador e denominador:
expanda $X(z)/z$ e use
$A_i=\left.(1-p_iz^{-1})X(z)\right|_{z=p_i}$.

\paragraph{Atraso e implementação.} $z^{-1}$ é o \textbf{operador de atraso
unitário}: cada $z^{-1}$ no numerador/denominador é uma amostra de memória. É
essa a ponte entre $H(z)$ e o fluxograma de implementação (Capítulo~\ref{sec:estr}).

\subsection*{Exercícios resolvidos}

% ---------------------------------------------------------------- 9.1
\exer{Calcule $X(z)$ e a ROC de $x[n]=(1/2)^nu[n]$; localize polos e zeros.}

\textbf{Solução.} Soma geométrica:
\[
  X(z)=\sum_{n\ge0}\left(\tfrac12 z^{-1}\right)^{n}
      =\frac{1}{1-\tfrac12z^{-1}}=\frac{z}{z-\tfrac12},\qquad
  |z|>\tfrac12 .
\]
Polo em $z=1/2$; zero em $z=0$. Como a ROC contém $|z|=1$, o sistema é
estável. Verificação: $\sum_n|h[n]|=\sum(1/2)^n=2<\infty$.

\resp{$X(z)=\dfrac{z}{z-1/2}$, ROC $|z|>1/2$, polo em $0{,}5$ e zero na
origem.}

% ---------------------------------------------------------------- 9.2
\exer{Determine $x[n]$ (causal) para $X(z)=\dfrac{z^2}{(z-1/2)(z-1/4)}$.}

\textbf{Solução.} Dividindo por $z$:
\[
  \frac{X(z)}{z}=\frac{z}{(z-1/2)(z-1/4)}
  =\frac{A}{z-1/2}+\frac{B}{z-1/4},
\]
\[
  A=\frac{1/2}{1/2-1/4}=2,\qquad
  B=\frac{1/4}{1/4-1/2}=-1 .
\]
Logo $X(z)=\dfrac{2z}{z-1/2}-\dfrac{z}{z-1/4}$ e, para a ROC $|z|>1/2$
(causal),
\[
  x[n]=\left[2\left(\tfrac12\right)^{n}-\left(\tfrac14\right)^{n}\right]u[n].
\]
Valores: $x[0..5]=\{1;\ 0{,}75;\ 0{,}4375;\ 0{,}234375;\ 0{,}12109375;\
0{,}06152344\}$ — checagem imediata: $x[0]=1$ porque o numerador e o
denominador têm a mesma ordem em $z$.

\resp{$x[n]=\left[2(1/2)^n-(1/4)^n\right]u[n]$, com $x[0]=1$ e $x[1]=0{,}75$.}

% ---------------------------------------------------------------- 9.3
\exer{Para $X(z)=\dfrac{1}{(1-\frac12z^{-1})(1-2z^{-1})}$, discuta as três
possíveis ROCs quanto a causalidade e estabilidade.}

\textbf{Solução.} Polos em $z=1/2$ e $z=2$. As ROCs possíveis (anéis
delimitados por polos) são:

\begin{center}
\small
\begin{tabular}{llll}
\toprule
ROC & Sequência & Causal? & Estável?\\
\midrule
$|z|>2$ & $x[n]=\bigl[-\frac13(\frac12)^n+\frac43\,2^{n}\bigr]u[n]$
        & sim & \textbf{não} (fora $|z|=1$)\\
$\frac12<|z|<2$ & $x[n]=-\frac13(\frac12)^nu[n]-\frac43\,2^nu[-n-1]$
        & \textbf{não} & \textbf{sim} (contém $|z|=1$)\\
$|z|<\frac12$ & $x[n]=\bigl[\frac13(\frac12)^n+\frac43\,2^n\bigr]u[-n-1]$ & não & não\\
\bottomrule
\end{tabular}
\end{center}

Decomposição em frações parciais (usada na tabela):
\[
  X(z)=\frac{-1/3}{1-\frac12z^{-1}}+\frac{4/3}{1-2z^{-1}},
  \qquad
  \text{checagem: } X(1)=-\tfrac13\cdot2+\tfrac43\cdot(-1)=-2 \ \checkmark
\]

\resp{A única ROC estável é $\frac12<|z|<2$, que corresponde a um sistema
\emph{não-causal}. Querer causal e estável ao mesmo tempo seria impossível —
é exatamente a situação do canal $H(z)=1-2z^{-1}$ do Desafio 04.}
